\chapter*{Leituras Essenciais}
\addcontentsline{toc}{chapter}{Leituras Essenciais}
%================================================

\section*{Como utilizar esta bibliografia}

A literatura sobre Modelagem de Distribuição de Espécies (Species Distribution
Modeling -- SDM) cresceu exponencialmente nas últimas décadas, acompanhando o
avanço da ecologia quantitativa, da biogeografia, da ciência de dados
ambientais e das mudanças climáticas.

Embora milhares de artigos tenham sido publicados sobre o tema, algumas obras
se destacam por estabelecer os fundamentos conceituais, metodológicos e
computacionais que estruturam a área atualmente.

Esta seleção reúne livros, artigos clássicos e trabalhos de referência que
servem como ponto de partida para estudantes, pesquisadores e profissionais
interessados em aprofundar seus conhecimentos em modelagem de nicho ecológico,
distribuição de espécies, mudanças climáticas e conservação da biodiversidade.

\section{Livros fundamentais}

\begin{table}[H]
	\centering
	\caption{Principais livros recomendados para aprofundamento em SDM.}
	\label{tab:livros_sdm}
	
	\small
	
	\begin{tabular}{
			p{0.30\textwidth}
			p{0.25\textwidth}
			p{0.35\textwidth}
		}
		
		\toprule
		
		\textbf{Livro} &
		\textbf{Autores} &
		\textbf{Contribuição}
		
		\\
		\midrule
		
		\textit{Mapping Species Distributions} &
		Franklin (2010) &
		Referência clássica sobre fundamentos ecológicos, estatísticos e espaciais da SDM.
		
		\\
		
		\textit{Habitat Suitability and Distribution Models} &
		Guisan, Thuiller \& Zimmermann (2017) &
		Obra moderna e abrangente sobre teoria, algoritmos e aplicações.
		
		\\
		
		\textit{Ecological Niches and Geographic Distributions} &
		Peterson et al. (2011) &
		Fundamentos conceituais do nicho ecológico e modelagem biogeográfica.
		
		\\
		
		\textit{Species Distribution Modeling with R} &
		Hijmans \& Elith &
		Aplicações práticas em ambiente R.
		
		\\
		
		\textit{Biogeography} &
		Lomolino et al. (2017) &
		Fundamentos biogeográficos necessários para interpretação dos modelos.
		
		\\
		
		\bottomrule
		
	\end{tabular}
	
\end{table}

\section{Artigos clássicos e fundamentais}

Os trabalhos apresentados a seguir são amplamente reconhecidos como marcos
conceituais e metodológicos da Modelagem de Distribuição de Espécies.

\begin{table}[H]
	\centering
	\caption{Artigos fundamentais para compreensão da evolução da SDM.}
	\label{tab:artigos_classicos_sdm}
	
	\small
	
	\begin{tabular}{
			p{0.30\textwidth}
			p{0.18\textwidth}
			p{0.40\textwidth}
		}
		
		\toprule
		
		\textbf{Referência} &
		\textbf{Ano} &
		\textbf{Contribuição}
		
		\\
		\midrule
		
		Guisan \& Zimmermann &
		2000 &
		Síntese pioneira sobre modelos preditivos de distribuição de espécies.
		
		\\
		
		Phillips, Anderson \& Schapire &
		2006 &
		Introdução formal do algoritmo Maxent.
		
		\\
		
		Elith et al. &
		2006 &
		Comparação abrangente de métodos modernos de SDM.
		
		\\
		
		Araújo \& New &
		2007 &
		Fundamentos da modelagem ensemble.
		
		\\
		
		Soberón \& Peterson &
		2005 &
		Integração entre nicho ecológico e distribuição geográfica.
		
		\\
		
		Barve et al. &
		2011 &
		Importância da definição da área acessível ($M$).
		
		\\
		
		Elith et al. &
		2010 &
		Avaliação da extrapolação ambiental.
		
		\\
		
		Zurell et al. &
		2020 &
		Protocolo internacional para relatórios e reprodutibilidade em SDM.
		
		\\
		
		\bottomrule
		
	\end{tabular}
	
\end{table}

\section{Leituras sobre mudanças climáticas e biodiversidade}

A integração entre SDM e mudanças climáticas tornou-se uma das principais áreas
de aplicação da modelagem ecológica. As obras abaixo são particularmente
relevantes para estudos de projeções futuras e conservação.

\begin{itemize}
	
	\item Thuiller et al. (2005) --- Climate change threats to plant diversity in Europe.
	
	\item Araújo et al. (2019) --- Standards for distribution models in climate change studies.
	
	\item IPCC (2021) --- Climate Change 2021: The Physical Science Basis.
	
	\item IPBES (2019) --- Global Assessment Report on Biodiversity and Ecosystem Services.
	
	\item Pecl et al. (2017) --- Biodiversity redistribution under climate change.
	
\end{itemize}

\section{Leituras sobre paleoclima e estabilidade climática}

A reconstrução de distribuições passadas representa uma importante aplicação da
SDM para investigação de refúgios climáticos, estabilidade ambiental e
história biogeográfica das espécies.

\begin{itemize}
	
	\item Svenning et al. (2011).
	
	\item Carnaval \& Moritz (2008).
	
	\item Graham et al. (2010).
	
	\item Peterson et al. (2011).
	
	\item Nogués-Bravo (2009).
	
\end{itemize}

\section{Leituras sobre reprodutibilidade científica}

A transparência e a reprodutibilidade tornaram-se componentes centrais da
pesquisa ecológica moderna. As referências abaixo fornecem excelentes pontos de
partida para implementação de boas práticas em projetos de SDM.

\begin{itemize}
	
	\item Peng (2011) --- Reproducible Research in Computational Science.
	
	\item Sandve et al. (2013) --- Ten Simple Rules for Reproducible Computational Research.
	
	\item Wilkinson et al. (2016) --- FAIR Guiding Principles.
	
	\item Zurell et al. (2020) --- Standard protocol for reporting species distribution models.
	
\end{itemize}

\section{Sequência recomendada de estudo}

Para leitores que desejam aprofundar seus conhecimentos em SDM, recomenda-se a
seguinte sequência de leitura:

\begin{enumerate}
	
	\item Guisan \& Zimmermann (2000);
	
	\item Franklin (2010);
	
	\item Peterson et al. (2011);
	
	\item Elith et al. (2006);
	
	\item Barve et al. (2011);
	
	\item Guisan, Thuiller \& Zimmermann (2017);
	
	\item Zurell et al. (2020).
	
\end{enumerate}

Essa sequência conduz o leitor desde os fundamentos ecológicos e
biogeográficos da modelagem até as práticas modernas de validação,
reprodutibilidade e aplicação em conservação sob mudanças climáticas.

\begin{dica}
	\textbf{Mensagem final ao leitor.}
	
	A Modelagem de Distribuição de Espécies é uma área em rápida evolução.
	Nenhum livro ou curso consegue esgotar completamente os conceitos,
	algoritmos e aplicações disponíveis.
	
	O desenvolvimento contínuo depende da leitura crítica da literatura
	científica, da prática com dados reais e da integração entre teoria
	ecológica, análise quantitativa e conservação da biodiversidade.
\end{dica}

